\section{Discussion}
\label{sec:discussion}

\subsection{Why a few days is enough}
\label{sec:discussion-fewdays}

The paper's central practical result is that EMOS pooled crosses
TimesFM-3's CRPS at \BpCrpsPooledCases{} cases (\BpCrpsPooledDays{}
calendar days, Section~\ref{sec:results-breakpoint}) -- but that raw
crossing is a point estimate, not yet a statistically defensible claim.
Testing significance across the low-\(N\) range shows the earliest
training size at which the advantage is significant under both a
day-blocked \(t\)-test and a Wilcoxon test is \BpCrpsPooledSigK{} cases
(\BpCrpsPooledSigDays{} days, Section~\ref{sec:results-significance}) --
day-blocking being our primary block definition; the station-blocked
\(t\)-test never reaches significance in the tested \(k=1\)--9 range,
while the station-blocked Wilcoxon test does from
\(k=\BpCrpsPooledStationWilcoxonFirstSigK{}\) onward, the same
divergence already reported at \SigTestNCases{} cases now shown to
persist throughout -- still a small absolute training size, and still
well inside the regime
\citet{baran_2016_combining} and \citet{gebetsberger_2018_estimation}
already identify as usable for station-adaptive EMOS, and that the
time-series and lead-time-continuous EMOS variants of
\citet{jobst_2024_ts_emos} and \citet{wessel_2024_leadtime_continuous}
target directly -- but nearly twice the raw crossing's calendar-day
figure, not the same number. The practical implication for a newly
commissioned station is direct, stated at the number the evidence
actually supports: roughly \BpCrpsPooledSigDays{} days of paired data,
not months, are enough to stop relying on a zero-shot foundation model's
raw CRPS with statistical confidence; the raw crossing at
\BpCrpsPooledDays{} days is the earliest the two methods' CRPS point
estimates change order, not the earliest that difference is established
rather than noise.

This result is deliberately scoped to CRPS alone. Coverage tells a
different story at substantially more data -- \CoverageCrpsRatioPooled{}x
later for pooled EMOS, \CoverageCrpsRatioLocal{}x later for local EMOS
(ratios from the unrounded breakpoint estimates, not the values as
displayed; Section~\ref{sec:results-breakpoint}: \BpCoveragePooledCases{} cases
pooled, \BpCoverageLocalCases{} cases locally), and interval width has no
single crossing point at all -- EMOS pooled starts sharper than TimesFM-3
and crosses to wider as it is given more data, while EMOS local stays
sharper throughout the tested range (Section~\ref{sec:breakpoint-def}'s
three-outcome distinction exists precisely so this kind of metric-specific
divergence is reported as what it is, not flattened into one number). "A
few days is enough" is therefore true for CRPS specifically, not for every
metric this paper reports -- Sections~\ref{sec:discussion-trivial}
and~\ref{sec:discussion-mechanism} address what CRPS alone does and does
not establish about why.

The question this paper asks -- how much training data a trained method
needs to replace a zero-shot foundation model -- has a deployment-cost
dimension this paper has not systematically measured: TimesFM-3 requires
a GPU forward pass through a 330-million-parameter pretrained model for
every inference call, while fitting EMOS end-to-end on the full training
archive took \EmosPooledFitSeconds{}~s (pooled) or
\EmosLocalFitSeconds{}~s (local) of CPU time in this study's own runs, and
inference afterward is a closed-form Gaussian evaluation, not a further
optimization. Once the training-data threshold in this paper is crossed,
the practical choice is not a single-axis trade-off between accuracy and
convenience: the trained method is both more accurate and cheaper to run.

\subsection{Positioning against trivial baselines}
\label{sec:discussion-trivial}

Section~\ref{sec:results-full}'s full-data comparison (Table~\ref{tab:full-comparison})
shows a result that splits by metric in a way worth stating plainly rather
than softening: on CRPS, a one-parameter variance-inflation rescaling of
the raw ensemble already beats TimesFM-3, and the fixed-climatological
variant (b, \CrpsInflBFull{} K vs.\ TimesFM-3's \CrpsTsfmFull{} K) does so
with no training data at all -- squarely the pattern
\citet{zhang_2025_context_parroting} describes for trivial baselines
against zero-shot foundation models generally, now confirmed in station
postprocessing specifically. A comparable dynamic -- a pretrained model
matching, but not decisively beating, a much simpler well-tuned
alternative -- is reported in a different domain by
\citet{molinaro_2026_universal_downscaling}.

On the calibration-sharpness axis, the evidence points the other way,
and this finding is now settled rather than conditional: the
coverage-matched variant (c), its multiplier fit on training data alone
(never on test-set coverage, Section~\ref{sec:methods}), needs
\WidthInflCFull{} K to reach as close to TimesFM-3's own coverage as a
training-fit multiplier allows -- substantially wider than TimesFM-3's
\WidthTsfmFull{} K. Variant (c) is leakage-free by design, which means it
is not an upper bound on what a trivial one-parameter rescaling could
achieve at TimesFM-3's coverage target -- only a lower bound on how well
a \emph{legitimately deployable} one can do. The leakage-optimistic
counterpart -- a multiplier fit directly on the test set's own coverage,
answering "what is the best any trivial rescaling could ever do here,
even with an unfair advantage" -- needs \InflLeakyWidth{} K, still wider
than TimesFM-3's \WidthTsfmFull{} K, at matched coverage: the leaky
variant's root-finder converges to \InflLeakyCoverage{}, identical to
TimesFM-3's \CoverageTsfmFull{} target to within floating-point
precision, so the width comparison is not confounded by a residual
coverage gap the way variant (c)'s is. This margin is not a bare point
estimate: a day-blocked bootstrap on the per-instance width difference
gives \InflLeakyWidthDiffPoint{} K (95\% CI
[\InflLeakyWidthDiffCiLow{}, \InflLeakyWidthDiffCiHigh{}]) -- entirely
positive, so the difference is statistically distinguishable from zero,
not merely numerically positive on average. The interval is conditional
on the leakage-optimistic multiplier fit once on the full test set, not
re-fit within each bootstrap resample: it quantifies sampling variability
in the width difference at that fixed multiplier, not variability in the
multiplier itself. This margin is also not explained by EMOS's own
training/test ensemble-size mismatch: subsampling the test side down to
the training side's member count changes EMOS's own interval width by no
more than \EnsSubsampleWidthDiffLocal{} K (Section~\ref{sec:limitations}),
under one-tenth of this margin. TimesFM-3's
calibration-sharpness advantage is therefore real: it survives even this
unrealistically favorable trivial baseline, with uncertainty quantified
(subject to that conditioning), so \citet{zhang_2025_context_parroting}'s
critique, while correct on CRPS (above), does not extend to this axis.

The honest summary of Section~\ref{sec:results-full} is: TimesFM-3 loses
to a trivial rescaling on CRPS, unconditionally; it beats even a
leakage-optimistic trivial rescaling on sharpness at matched coverage,
so the calibration-sharpness advantage is not a leakage-free-vs.-leaky
comparison artifact.

\subsection{The short-lead advantage: information asymmetry or architecture}
\label{sec:discussion-mechanism}

\textbf{This section remains provisional; the architectural-advantage
hypothesis has not been confirmed.} Section~\ref{sec:covariates}
proposed a single mechanism for TimesFM-3's short-lead CRPS advantage
(Section~\ref{sec:results-leadtime}: first sign flip at
\FirstSignFlipHours{}~h, durable crossover at
\DurableCrossoverHours{}~h): access to two channels EMOS and DRN were not
given in this study's configuration (Table~\ref{tab:info-channels}) --
the station's own recent observation history, and ensemble spread. These
are two separate questions with two separate tests; one has now been
run, with a limited, inconclusive result, and one is still outstanding.

\textbf{Observation history (tested; does not support the
information-asymmetry explanation, but does not settle the question
either).} A persistence-augmented EMOS at 0--24~h tests whether adding
the station's own last observation as a predictor closes the CRPS gap
toward TimesFM-3. It does not, substantially: only
\PersistenceGapClosedPct{}\% of the gap closes (persistence weight
\PersistenceWeight{}, matched subset, \NPersistenceMatchedInstances{}
instances), confirmed after ruling out lead-0 contamination as a
confound (Section~\ref{sec:limitations}: exact-match and trend checks
both clean). This result does not support the information-asymmetry
explanation as stated -- the simplest way of giving EMOS the missing
observation-history channel does not close the gap. It does not, on its
own, establish that the advantage is architectural instead: the test has
three real limitations that keep it from distinguishing the two
hypotheses cleanly, not just describing them as a formality.

\begin{itemize}
\item \emph{Lead-pooled fit.} EMOS pools all \NLeadTimes{} lead times
into one fit (Section~\ref{sec:methods-list}); the persistence term is
fit once, lead-blind, inside that pooled model, while TimesFM-3's actual
covariate and target information vary naturally with lead time. The two
methods are not given the same structural opportunity to use
lead-dependent information -- this is not an equal comparison.
\item \emph{One linear term, one lag.} The persistence predictor is a
single linear term on the single most recent observation. TimesFM-3
receives a context window of 40 or more past observations, processed by
a sequence model that can learn arbitrarily nonlinear, multi-lag
structure. Adding one linear lag does not give EMOS "comparable
information" to a 40-observation context in any meaningful sense --
only a much smaller, cruder approximation to it.
\item \emph{Metric-dependent.} In the same 0--24~h bucket, TimesFM-3
wins on CRPS but EMOS is significantly \emph{closer to nominal coverage}
(day-blocked \(t\): \(p=\LeadSigCoverageBucketShortDayT\), EMOS's favor,
Section~\ref{sec:results-significance}). Whatever "architectural
advantage" this test might support holds for CRPS specifically, not
uniformly across every metric in the same lead-time bucket.
\end{itemize}

Given these three limitations, the honest reading of this test is
narrower than "architecture, not information asymmetry": it shows that
the simplest information-asymmetry fix is insufficient, not that
architecture is confirmed as the explanation. A test able to distinguish
the two hypotheses would need a lead-resolved EMOS fit (not pooled) with
a multi-lag persistence predictor (not one linear term on one lag) --
this paper does not run that test, and does not claim the weaker test
run here settles the question.

\textbf{Ensemble spread / calibration (resolved: does not support the
asymmetry framing).} The other half of the proposed asymmetry -- that
TimesFM-3's calibration deficit relative to EMOS reflects a missing
spread channel, not an architectural limit -- has now been tested
directly. The model's covariate channel does accept spread as a second,
simultaneous input, and the model is genuinely sensitive to that
channel's real informational content, not merely its presence (a
preliminary test confirms both, Section~\ref{sec:covariates},
Section~\ref{sec:limitations}), but supplying it this way did not close
the calibration deficit: CRPS and interval width both worsen, and
coverage moves only marginally closer to nominal at a substantially
wider interval (Section~\ref{sec:discussion-covariate-check} reports the
full result). This does not confirm the asymmetry framing for this half
of the mechanism question, and the observation-history result above does
not settle it either: the two channels test genuinely different things,
and a finding for one is not evidence for the other.

\subsection{The two-covariate evaluation: does spread change calibration?}
\label{sec:discussion-covariate-check}

Section~\ref{sec:covariates}'s preliminary test (T1) established that
TimesFM-3's covariate channel accepts ensemble spread as a second,
simultaneous input alongside the ensemble mean, and that the model's
output is genuinely sensitive to that channel's real informational
content. This section reports the full-scale evaluation that
preliminary result motivated: supplying both the ensemble mean and
ensemble spread as simultaneous past-future covariates, on the same
\NMatchedInstances{} matched instances used for every other TimesFM-3
number in this paper, and comparing the result against the mean-only
configuration used everywhere else.

Supplying spread did not improve TimesFM-3's CRPS: it worsens from
\CrpsTsfmFull{} K to \CrpsTsfmTwoCov{} K, a net negative of
\CrpsTwoCovDiff{} K. No uncertainty interval has been computed for this
difference; its magnitude is comparable to the ensemble-size sensitivity
measured directly in Section~\ref{sec:limitations}
(\EnsSubsampleCrpsDiffPooled{}--\EnsSubsampleCrpsDiffLocal{} K), so it
should be read as a real but not dramatically large effect, not a
precisely bounded one. Coverage moves closer to the 0.80-nominal target,
from \CoverageTsfmFull{} to \CoverageTsfmTwoCov{}: in practice, the
two-covariate configuration is calibrated at this target, where the
mean-only configuration under-covers. Interval width also increases,
from \WidthTsfmFull{} K to \WidthTsfmTwoCov{} K (\WidthTwoCovDiff{} K
wider) -- but by this paper's own standard of judging sharpness jointly
with calibration (Section~\ref{sec:methods},
\citealp{gneiting_2007_calibration_sharpness}), these two widths are not
directly comparable on the same axis: they are measured at different
coverage levels, and a wider interval is the expected cost of a
configuration that is closer to properly calibrated, not on its own
evidence of a worse model. This reading has independent support: fitting
the leakage-optimistic multiplier (Section~\ref{sec:discussion-trivial})
at this same \CoverageTsfmTwoCov{} target, rather than TimesFM-3's
original \CoverageTsfmFull{}, still needs \LeakyWidthAtTwoCovCoverage{}
K -- \LeakyWidthAtTwoCovCoverageDiff{} K wider than the two-covariate
configuration's own \WidthTsfmTwoCov{} K at the identical coverage.
Finding 2's width advantage therefore holds at two independent coverage
points, not only at TimesFM-3's original one.

Supplying ensemble spread as a second covariate therefore worsens
TimesFM-3's CRPS and moves it along the calibration-sharpness trade-off
-- closer to nominal coverage, at a wider interval -- rather than
cleanly closing the calibration deficit: this is a shift along that
trade-off, not a clean improvement or a clean failure. This result does
not confirm the information-asymmetry framing for the calibration half
of the mechanism question (Section~\ref{sec:covariates}): TimesFM-3 was
not simply prevented by an input-design choice from recovering EMOS-like
calibration once that choice was corrected. But it does not confirm
architecture as the explanation instead, either: the model accepted the
second input and was demonstrably sensitive to it
(Section~\ref{sec:limitations}), and did shift its output along the
calibration-sharpness trade-off when given it -- it just did not convert
that shift into a net improvement on CRPS, or into a clean win on the
paper's joint sharpness-calibration standard, within this study's
zero-shot, frozen-weights setup. Whether a fine-tuned or differently
prompted configuration would do better remains untested and outside this
paper's scope; so does whether a different input scale or normalization
convention for the spread channel would change this result
(Section~\ref{sec:limitations}).

\subsection{Implementation pitfall: covariate wiring}
\label{sec:pitfall}

\citet{kreusel_2026_covariates_key} reports that covariate handling, not
architecture or scale, is the dominant factor in time series foundation
model performance on real tasks. This project's own development history is
a direct, first-hand instance of that finding rather than a citation of
it: an earlier version of TimesFM-3's covariate pipeline
(Section~\ref{sec:covariates}) supplied an incoherent covariate channel --
one that ran without error and produced plausible-looking quantile output
-- and measurably degraded TimesFM-3's apparent skill before the bug was
found and corrected. The pipeline did not crash; it silently produced a
worse model. The practical warning for anyone implementing a
covariate-aware time series foundation model pipeline is that "it runs and
returns sensible-looking numbers" is not evidence the covariate is wired
correctly -- only a dedicated correctness check on the covariate channel
itself (not just on the model's output shape or its overall CRPS) would
have caught this before it silently degraded results. This experience is
the direct reason this project later required a distinguishing
correctness check -- not just "does the call succeed," but "is the
output genuinely sensitive to this covariate's real informational
content, not merely its presence" -- before trusting any further
covariate-channel change (Section~\ref{sec:discussion-covariate-check}'s
ablation test is that check, run precisely because of the lesson learned
here).

\subsection{The value of seasonal coverage in small training sets}
\label{sec:discussion-seasonal}

Section~\ref{sec:data}'s two sampling arms answer different questions by
construction: the contiguous arm's small-\(k\) subsets see less seasonal
variety than a full archive would, while the random arm's uniform draw
gives artificial full-season coverage even at small \(k\). At
\SeasonalCompareNDays{} days, EMOS pooled's coverage is already visibly
different between the two arms -- \CoverageRandomSeasonalN{} under the
random arm against \CoverageContigSeasonalN{} under the contiguous arm --
while CRPS is nearly identical (\CrpsRandomSeasonalN{} K vs.\
\CrpsContigSeasonalN{} K). This is consistent with
Section~\ref{sec:discussion-fewdays}'s CRPS-crosses-fast,
coverage-crosses-slow pattern having a seasonal-coverage component: a
model can already predict the mean well from a handful of contiguous days,
but calibrating the predictive interval correctly needs to have seen a
representative spread of conditions, not just recent ones. Operationally,
this means the realistic deployment scenario -- a newly commissioned
station with only contiguous recent history, the contiguous arm's premise
-- reaches good CRPS quickly but should not be assumed to be well
calibrated at the same training size a full-season random sample would be.
\citet{baran_2026_sharpness_penalty} targets the same sharpness axis on
the same benchmark; the overlap is in the axis studied, not in this
specific contiguous-vs-random comparison, which does not appear in that
work.

\subsection{Limitations}
\label{sec:limitations}

\textbf{Covariate channel design; tested, but does not resolve the
mechanism question.} TimesFM-3's covariate channel was given only the
ensemble mean in this study's main comparison
(Section~\ref{sec:covariates}); whether supplying ensemble spread too
would change calibration was tested directly
(Section~\ref{sec:discussion-covariate-check}). The channel does accept
a second, simultaneous input, and the model is genuinely sensitive to
its real informational content, not merely its presence -- but CRPS
worsens when spread is supplied this way, and while coverage moves
closer to nominal (a genuine calibration shift, not merely a wider
interval for its own sake), the interval widens substantially to get
there, so the net result is a shift along the calibration-sharpness
trade-off, not a clean win (Section~\ref{sec:discussion-covariate-check}).
This result does not confirm the information-asymmetry framing (spread
was not simply withheld from a channel that would otherwise have used it
straightforwardly beneficially), but it does not confirm architecture as
the explanation either: the channel accepted the input and the model was
demonstrably sensitive to it, and did shift its behavior in response to
it, without that shift converting into a net improvement. Whether a
fine-tuned or differently prompted configuration would use spread more
successfully is untested and outside this paper's scope. We report this
as an unresolved mechanism question rather than resolving it either way
beyond what this evidence supports.

\textbf{Spread channel scale and normalization.} Ensemble spread was
supplied to TimesFM-3's second covariate channel as a standard
deviation, in the same physical units (Kelvin) as the ensemble mean but
roughly two orders of magnitude smaller in raw value (typically 1--3~K
vs.\ 270--290~K) -- a real scale disparity between the two channels as
supplied. This is not, on its own, an unaddressed confound: TimesFM-3's
own preprocessing z-normalizes each covariate channel independently,
computing a running mean and standard deviation per channel rather than
sharing statistics across channels (confirmed directly from the
installed \texttt{timesfm3} package's source, not inferred), so the raw
magnitude difference between the mean and spread channels does not by
itself explain why supplying spread failed to improve performance. The
spread channel was supplied at a single scale and as a single channel,
per-channel normalization is confirmed rather than merely assumed, and
the result above should be read as "this specific input design does not
help," not as "TimesFM-3 cannot use ensemble spread" more broadly --
untested alternatives include a differently transformed spread channel
(e.g.\ log-spread) or a channel expressing spread relative to the mean.

\textbf{Lead-time pooling.} EMOS and DRN each fit a single parameter set
pooled across all \NLeadTimes{} lead times (Section~\ref{sec:methods-list}),
while TimesFM-3 and the raw ensemble vary naturally by lead time. A
lead-blind fit is necessarily too wide at short lead (where true
uncertainty is small) and too narrow at long lead (where it is large)
relative to a lead-time-aware fit, which means the CRPS crossings reported
in Section~\ref{sec:results-breakpoint} are, if anything, understated for
EMOS and DRN, not artifacts of this pooling. It is a genuine limitation
for the sharpness comparison specifically, since a lead-blind interval's
width behavior is not directly comparable to TimesFM-3's per-lead-varying
output. For DRN in particular, this feature-constrained setup is not the
architecture's originally-intended one (\citealp{rasp_2018_drn} additionally
use auxiliary NWP fields and lead time as inputs): where DRN underperforms
in this paper, the correct reading is that this architecture, under this
feature-constrained setup, was outperformed -- not that the architecture
itself is weak.

\textbf{Persistence test run on a lead-pooled EMOS fit.} The
persistence-augmented EMOS test (Section~\ref{sec:discussion-mechanism})
inherits this same lead-pooling limitation directly: the persistence
predictor is fit once, lead-blind, inside a model already disadvantaged
by pooling across all 21 lead times against a naturally lead-varying
TimesFM-3. This reduces the test's power to distinguish an
information-asymmetry explanation from an architectural one -- a
lead-resolved EMOS fit with a multi-lag persistence predictor would be a
meaningfully stronger test, and this paper does not run one.

\textbf{Target-variable specificity.} The breakpoint measured here is
for \(2\,\mathrm{m}\) temperature, where EMOS's Gaussian predictive
family is close to ideally specified. For targets requiring non-Gaussian
or target-specific reformulations (Section~\ref{sec:intro}: e.g.\
\citealp{scheuerer_2013_precip_emos}'s precipitation-specific EMOS), the
trained method's small-sample advantage would not necessarily appear as
early.

\textbf{Tail truncation in twCRPS.} The tail-weighted CRPS defined in
Section~\ref{sec:metrics} uses the upper two levels of a fixed
\NQuantileLevels{}-level quantile grid (0.8, 0.9) as its tail weighting.
This tests the upper end of the quantile grid this study fixed for every
method, not whether TimesFM-3 genuinely lacks tail behavior beyond p90 --
the grid itself has no quantile levels past 0.9 for any method, so no
comparison in this paper can distinguish "no tail" from "a tail this
grid cannot see."

\textbf{PIT diagnostics are descriptive only.}
\citet{broecker_2018_serial_dependence} and
\citet{broecker_2020_stratified_rank} show that formal rank-uniformity
hypothesis tests require correction for serial dependence, which this
study's PIT values genuinely exhibit (station~$\times$~time~$\times$~lead
panel structure, Section~\ref{sec:significance}); no such test is applied
to the PIT histograms in Section~\ref{sec:results-sharpness}, which are
read descriptively for this reason.
\citet{dirkson_2025_misdiagnosing_reliability} further shows that coverage
and rank-based reliability diagnostics can both be misled by the same
covariance pathology -- we do not claim the descriptive PIT reporting in
this paper is immune to that critique; it is a genuine, unresolved
limitation, not one worked around by reporting PIT histograms instead of
a formal test.

\textbf{Random-arm seed count.} The random sampling arm's mean and
standard deviation (Section~\ref{sec:data}) are estimated from
\DataSizeSweepSeeds{} seeds. This is enough to see the arm's qualitative
pattern relative to the contiguous arm, but the standard deviation itself
is estimated noisily from so few draws -- it should be read as an
indicator of spread, not a precise quantity, and Figure~\ref{fig:f1}'s
random-arm error bars inherit this same noise.

\textbf{Ensemble size mismatch.} EMOS's training-side ensemble variance is
estimated from \TrainEnsembleMembers{} members against
\TestEnsembleMembers{} on the test side (Section~\ref{sec:data}); an
11-member variance estimate is intrinsically noisier than a 51-member one
of the same true spread, so EMOS is trained against a noisier variance
estimate than the one used to evaluate it. This is no longer an
unbounded concern: subsampling 11 of the test side's 51 members (10
draws, recomputing the ensemble mean/variance from just those 11 and
re-evaluating the already-fitted full-\(N\) EMOS models) changes pooled
EMOS's own CRPS by \EnsSubsampleCrpsDiffPooled{} K, its interval width by
\EnsSubsampleWidthDiffPooled{} K, and its coverage@80\% by
\EnsSubsampleCoverageDiffPooled{} (local EMOS: \EnsSubsampleCrpsDiffLocal{} K,
\EnsSubsampleWidthDiffLocal{} K, \EnsSubsampleCoverageDiffLocal{}
respectively) -- both width differences are under one-tenth of
Section~\ref{sec:discussion-trivial}'s \InflLeakyWidthDiffPoint{} K
leakage-optimistic-width margin. The ensemble-size mismatch is real and
measurable, but it is too small to be a material contributor to that
finding.

\textbf{Unverified chronological assumption.} The contiguous sampling
arm orders cases by descending \((\mathrm{year\_idx}, \mathrm{time\_idx})\),
assuming ascending \(\mathrm{year\_idx}\) is chronologically more recent
(Section~\ref{sec:data}). Only \(\mathrm{time\_idx}\)'s ordering was
independently verified against real decoded dates; \(\mathrm{year\_idx}\)'s
polarity was not. This bounds its impact to training sizes beyond the
209-date single-cycle boundary (\(k>209\), i.e.\ N greater than roughly
728 calendar days) -- below that boundary, every contiguous-arm result in
this paper is unaffected regardless of which polarity is correct.

\textbf{Training-data contamination risk.} This study does not verify that
TimesFM-3's pretraining corpus excludes the 2017--2018 test period used
here. Data contamination in time series foundation model evaluation is an
active, documented concern
(\citealp{li_2026_tsfmaudit,chiazagomekpere_2026_leakage_aware,balloccu_2024_leak_cheat_repeat}),
and benchmark-methodology work
(\citealp{aksu_2024_gift_eval,meyer_2025_rethinking_eval}) argues evaluation
protocols in this area have not caught up with the risk. We cannot rule
this out for TimesFM-3 on EUPPBench; a demonstrated case of the same risk
materializing in a different domain is given by
\citet{pan_2026_contamination_electricity}. If TimesFM-3's apparent
zero-shot skill benefits from contamination, its true zero-shot CRPS would
be worse than reported here, which would only strengthen this paper's
conclusion that trained methods overtake it quickly -- it does not
threaten the paper's central claim, but it is a real, unaddressed
limitation of the TimesFM-3 numbers taken in isolation.

\textbf{Lead-0 contamination, checked and ruled out.}
Section~\ref{sec:discussion-mechanism} depends on the 0--24~h bucket,
which includes \(\mathrm{step\_hours}=0\) -- the analysis time, where the
"forecast" could in principle be tightly coupled to the observation it
is verified against for every method, not only TimesFM-3. Two checks
rule this out: the exact-match fraction between the ensemble-mean
forecast and the observation at lead 0 is \LeadZeroExactMatchFraction{}
(no higher than the \LeadZeroMaxOtherExactMatchFraction{} seen at any other
lead, both consistent with floating-point noise, not literal copying);
and lead 0's residual standard deviation (\LeadZeroActualStd{} K) is in
line with what a linear trend fit to leads \(\geq 6\)~h predicts at
\(\mathrm{step\_hours}=0\) (\LeadZeroPredictedStd{} K, ratio
\LeadZeroStdRatio{}) -- if anything very slightly above the trend, not a
discontinuous tightening below it. Every lead-time-resolved result in
this paper that includes lead 0 (Section~\ref{sec:results-leadtime}) can
be read without this particular caveat.
